%% toc: P2 9 9.5,9.6,9.7
\chapter{The Rules of \texorpdfstring{\(\mathsf S\)}{S} and the Classical Syllogisms}
\label{sec:p02-w02c-p02-002-rules}

The calculus \(\mathsf S\) of Chapter~\ref{sec:p02-syllogistic-inference}
is a modern formal system. Some of its rules are compact forms of
traditional syllogisms; others have no classical counterpart and answer to
the empty-term semantics or to the needs of the completeness proof. This
chapter sorts the rules accordingly, following Pratt-Hartmann and Moss's own
commentary on their rule sheet, proves that the nonclassical axiom T cannot
be removed, records why the redundant rule D1 is kept, and shows why the
literal-set lemmas used in the completeness proof are restricted to
universal premises.

\emph{Inputs.} The languages, set semantics and formula identifications of
Definitions~\ref{def:p02-syllogistic-language},
\ref{def:p02-syllogistic-semantics}
and~\ref{def:p02-syllogistic-identifications}; direct derivations
(Definition~\ref{def:p02-syllogistic-derivation}) and the rule sheets of
\(\mathsf S\) and \(\mathsf S^\dagger\) in
Chapter~\ref{sec:p02-syllogistic-inference}; the closure \(V^\Phi\) and its
lemmas (Definition~\ref{def:p02-syllogistic-literal-closure},
Lemmas~\ref{thm:p02-syllogistic-closure}
and~\ref{thm:p02-syllogistic-extension}); and the fourteen syllogism rules
read in \(\mathcal S\) (Table~\ref{tab:p02-w02-p02-002-fourteen}).

\emph{Conclusions.} D1, D2, D3, B and X have classical counterparts; A, T
and I do not. Without T neither calculus derives anything from the empty premise
set, so neither is complete
(Lemma~\ref{thm:p02-w02c-p02-002-t-needed}). D1 is redundant but is kept
for ease of use and for its link to the relational system of the same
source. The one-element model of Lemma~\ref{thm:p02-syllogistic-extension}
exists only because \(\Phi\) is universal: for the premise set
\(\{\exists(p,q)\}\) the closure construction yields a nonempty, consistent,
closed literal set none of whose complete extensions defines a model
(Lemma~\ref{thm:p02-w02c-p02-002-existential-closure}).

\section{Classical and nonclassical rules}
\label{sec:p02-w02c-p02-002-classical}

Pratt-Hartmann and Moss relate their rules to the traditional moods as
follows. D1 is one schema for the traditional syllogisms Darii and Ferio
\cite[pp.~8--9]{PrattHartmannMoss2008SyllogisticV1}, which are rows SR3 and
SR4 of Table~\ref{tab:p02-w02-p02-002-fourteen}. D2 and B they describe as
plainly forms of classical syllogisms; in that table, B yields Barbara and
Celarent, and D2 yields Disamis and Bocardo. D3, which the source says it
has met earlier in a more general form, yields Baroco there. X is the classical principle
\emph{ex falso quodlibet}, by which a contradiction licenses any conclusion.
It is a rule of the direct system, not reductio: it discharges no
assumption. Rules A, T and I have no classical counterparts
\cite[p.~12]{PrattHartmannMoss2008SyllogisticV1}.

None of A, T and I occurs in the derivations of
Table~\ref{tab:p02-w02-p02-002-fourteen}, which use B, D1, D2 and D3 alone.
The three nonclassical rules serve the empty-term semantics and the
completeness argument.
\begin{itemize}
\item \emph{Rule A}, from \(\forall(p,\bar p)\) to \(\forall(p,l)\), rests on
 empty extensions. If every \(p\) is a non-\(p\), nothing is a \(p\), and
 then every \(p\) is an \(l\) vacuously
 \cite[p.~12]{PrattHartmannMoss2008SyllogisticV1}. The proof of
 Lemma~\ref{thm:p02-syllogistic-extension} uses A to exclude the last
 possible conflict when a literal set is completed.
\item \emph{Rule T}, the axiom \(\forall(p,p)\), cannot be dropped: without
 it there would be no derivation of \(\forall(p,p)\) from the empty premise
 set \cite[p.~12]{PrattHartmannMoss2008SyllogisticV1}.
 Lemma~\ref{thm:p02-w02c-p02-002-t-needed} proves this, and more, from the
 definition of a derivation. T also gives the inclusion
 \(V\subseteq V^\Phi\) of Lemma~\ref{thm:p02-syllogistic-closure}.
\item \emph{Rule I}, from \(\exists(p,l)\) to \(\exists(p,p)\), keeps a
 witness and forgets one of its properties. The completeness proof for
 \(\mathsf S^\dagger\) uses it when both literals of a requested existential
 conclusion follow from the same property of a given witness.
\end{itemize}

\begin{hblemma}{thm:p02-w02c-p02-002-t-needed}{proved=annex:p02-w02c-p02-002-t-needed}{Without T nothing is derivable from no premises}
Let \(\mathsf S^-\) be the rule set of \(\mathsf S\) with T deleted, and
\(\mathsf S^{\dagger-}\) that of \(\mathsf S^\dagger\) with T deleted.
\begin{enumerate}
\item No formula has a direct derivation from \(\varnothing\) using only
 rules of \(\mathsf S^-\), or only rules of \(\mathsf S^{\dagger-}\).
\item For every atom \(p\), \(\varnothing\models\forall(p,p)\), and
 \(\forall(p,p)\) belongs to both \(\mathcal S\) and \(\mathcal S^\dagger\).
 Hence neither \(\mathsf S^-\) nor \(\mathsf S^{\dagger-}\) is complete for
 its fragment.
\item With T, the one-node tree whose leaf is the T instance
 \(\forall(p,p)\) is a derivation of \(\forall(p,p)\) from \(\varnothing\)
 in \(\mathsf S\) and in \(\mathsf S^\dagger\).
\end{enumerate}
\end{hblemma}

The lemma concerns only derivations without premises. With premises, the
smaller rule sets still derive many formulas; for instance, B derives
\(\forall(p,o)\) from \(\forall(p,q)\) and \(\forall(q,o)\).

\begin{hbexample}{def:p02-w02c-p02-002-t-example}{A valid universal with no premises}
The sentence ``every oak is an oak'' is \(\forall(o,o)\). It is true in the
structure of Example~\ref{def:p02-syllogistic-empty-term}, where oak is
empty, and in every other structure. In \(\mathsf S\) its derivation is the
single T instance. If T is deleted, Lemma~\ref{thm:p02-w02c-p02-002-t-needed}
shows that no tree of rule instances can end in \(\forall(o,o)\) without a
premise, although no structure falsifies it.
\end{hbexample}

\section{Why the redundant rule D1 is kept}
\label{sec:p02-w02c-p02-002-d1}

D1 is redundant in \(\mathsf S\): after a formula identification, each of
its instances is an instance of D2 or of D3
(Lemma~\ref{thm:p02-syllogistic-rules}). The source keeps it for two stated
reasons. It is easier to use, and it makes the relation between
\(\mathsf S\) and the system \(\mathsf R\) of the source's Section~4.2 more
transparent \cite[pp.~12--13]{PrattHartmannMoss2008SyllogisticV1}.
\(\mathsf R\) is a finite rule set for the source's relational fragment,
which adds binary relations to \(\mathcal S\); the source shows it to be
sound and refutation-complete. The first rule of \(\mathsf R\) is D1 with the
literal \(l\) widened to a c-term
\cite[p.~20]{PrattHartmannMoss2008SyllogisticV1}. A c-term is either a
literal or a complex term built from an atom and a binary relation, read as a
noun phrase such as ``thing which admires some \(q\)''
\cite[p.~4]{PrattHartmannMoss2008SyllogisticV1}. For the fourteen rules,
keeping D1 means that Darii, Ferio, Festino, Datisi and Ferison are each
derived by one application of the same rule, after at most one formula
identification (Table~\ref{tab:p02-w02-p02-002-fourteen}).

\section{Universal premises in the closure lemmas}
\label{sec:p02-w02c-p02-002-universal}

The source states the closure \(V^\Phi\) and its three lemmas on literal
sets for a set \(\Phi\) of universal formulas
\cite[pp.~14--15]{PrattHartmannMoss2008SyllogisticV1}. The defining formula
of \(V^\Phi\) in Definition~\ref{def:p02-syllogistic-literal-closure} makes
sense for any premise set, since it mentions only BTA derivations. The next
lemma shows that the model-building half of
Lemma~\ref{thm:p02-syllogistic-extension} then fails.

\begin{hblemma}{thm:p02-w02c-p02-002-existential-closure}{proved=annex:p02-w02c-p02-002-existential-closure}{Closure with an existential premise}
Let \(p,q\) be distinct atoms and \(\Phi=\{\exists(p,q)\}\). Read
\(V^\Phi\), and \(\Phi\)-closed, by the formulas of
Definition~\ref{def:p02-syllogistic-literal-closure} with this \(\Phi\).
\begin{enumerate}
\item For literals \(l,m\),
 \(\Phi\vdash_{\mathrm{BTA}}\forall(l,m)\) if and only if \(l=m\).
\item \(V^\Phi=V\) for every set \(V\) of literals. In particular every
 literal set is \(\Phi\)-closed, and \(\{\bar p\}\) is nonempty, consistent
 and \(\Phi\)-closed.
\item \(\{\bar p\}\) has a consistent, complete, \(\Phi\)-closed extension.
 For every such extension \(W\), the one-element structure whose element
 satisfies exactly the literals in \(W\) falsifies \(\exists(p,q)\).
\end{enumerate}
\end{hblemma}

The completeness proofs never meet this case. They split a finite premise
set into its universal part \(\Phi\) and its existential part, apply the
closure only to the universal part, and treat the existential premises
through Lemma~\ref{thm:p02-syllogistic-separation} and the witness
literals they supply.

\index{syllogistic logic!classical and nonclassical rules}
\index{syllogistic logic!rule T}
\index{ex falso quodlibet}
